\scanpage{43}
\mspage{- IV-\handwritten{977}- \handwritten{12 -}}

$u:B\to A$, où $B=k\handwritten{[[}T_1,..,T_r\handwritten{]]}$ ($r$ entier quelconque) faisant de $A$ une $B$-algèbre finie, $\underline{m}$ est le seul point de $\operatorname{Spec}(A)$ au-dessus'\deleted{\unsure{du|de}}de l'idéal maximal \deleted{\illegible{1}} $\underline{n}$ de $B$ ; \deleted{\illegible{1}} le morphisme \deleted{\unsure{$\operatorname{Spec}(A)\to\operatorname{Spec}(B)$}} \inserted{\handwritten{$\operatorname{Spec}(A)\to\operatorname{Spec}(B)$}} ne peut être non ramifié en $\underline{m}$ que si $A/\underline{n}A$ est un corps, extension séparable de $k$ (17.4.1), ce qui implique que le corps résiduel $K$ de $A$ est une extension séparable de $k$. Si cette condition n'est pas remplie, la conclusion de (18.11.5\deleted{\unsure{p|b}}) ne peut jamais être vérifiée pour $\underline{p}=\underline{m}$, quel que soit l'entier $m$.

(ii) Supposons toujours $\underline{p}\ne\underline{m}$, mais considérons le cas où la condition de (18.11.5) relative à \inserted{\handwritten{l'existence de}} l'idéal minimal $\underline{q}$ n'est pas remplie. Soit $\underline{J}$ le noyau de l'homomorphisme canonique $A\to A_{\underline{p}}$, autrement dit l'ensemble des annulateurs des éléments de $A-\underline{p}$, et posons $A'=A/\underline{J}$ ; on a nécessairement $\underline{J}\subset\underline{p}$, et si l'on pose $\underline{p}'=\underline{p}/\underline{J}$, l'homomorphisme canonique $A_{\underline{p}}\to A'_{\underline{p}'}$ est \uline{bijectif}. En vertu de (18.11.3.\unsure{3|2}) et de (\uline{0},20.4.12), l'homomorphisme canonique \deleted{\unsure{$\Omega^1_{A/k}\otimes_A A'$}} $\widehat{\Omega}^1_{A/k}\otimes_A A'\to\widehat{\Omega}^1_{A'/k}$ est \uline{surjectif}, donc il en est de même de $(\widehat{\Omega}^1_{A/k})_{\underline{p}}\to(\widehat{\Omega}^1_{A'/k})_{\underline{p}'}$. D'autre part, l'homomorphisme canonique $A'\to A'_{\underline{p}'}$ est par définition \uline{injectif}, et les idéaux \inserted{premiers} minimaux de $A'$ correspondent donc biunivoquement à ceux de $A'_{\underline{p}'}$ ;
% slanted
\margin{\handwritten{\uline{\unsure{Vider|Viduer}} (ii) sauf \unsure{raison} \unsure{sérieuse} !}}
\deleted{\unsure{donc}} \inserted{\handwritten{par suite}} il y a un idéal minimal $\underline{q}'$ de $A'$ \deleted{\illegible{1}} contenu dans \deleted{\illegible{1}} $\underline{p}'$ et tel que $\dim(A'/\underline{q}')=\dim(A')$. On peut ainsi appliquer (18.11.5) à $A'$ et $\underline{p}'$, et il y a donc un $k$-homomorphisme local $u':B\to A'$, avec $B=k\handwritten{[[}T_1,..,T_m\handwritten{]]}$ tel que le morphisme correspondant \deleted{\unsure{$\operatorname{Spec}$}} \inserted{$f':$} $\operatorname{Spec}(A')\to\operatorname{Spec}(B)$ soit fini et non ramifié au point $\underline{p}'$. Mais on peut factoriser $u'$en $B\xrightarrow{u}A\xrightarrow{p}A'$ où \deleted{\illegible{3}} $p$ est l'homomorphisme canonique ; d'autre part, l'injection canonique $\operatorname{Spec}(A')\to\operatorname{Spec}(A)$ est un\deleted{e} \deleted{\unsure{homomorphisme}} \uline{isomorphisme local} au point $\underline{p}'$ en vertu de l'isomorphisme $A_{\underline{p}}\xrightarrow{\sim}A'_{\underline{p}'}$ et de (\uline{I},6.5.4). On voit donc que l'hypothèse a) de (18.11.5) entraîne encore, dans ce cas, l'existence d'un $k$-homomorphisme local $u:B\to A$ tel que

% archive sheet number: 42

